% Generated draft: requires mathematical and source review.
\documentclass[11pt]{article}
\usepackage[margin=25mm]{geometry}
\usepackage{fontspec}
\setmainfont{Latin Modern Roman}
\usepackage{amsmath,amssymb,mathtools,unicode-math}
\setmathfont{Latin Modern Math}
\usepackage{enumitem,needspace,xurl,hyperref}
\setlength{\parindent}{0pt}
\setlength{\parskip}{5pt}
\setlength{\emergencystretch}{3em}
\newcommand{\N}{\mathbb{N}}
\newcommand{\Z}{\mathbb{Z}}
\newcommand{\Q}{\mathbb{Q}}
\newcommand{\R}{\mathbb{R}}
\newcommand{\C}{\mathbb{C}}
\newcommand{\F}{\mathbb{F}}
\newcommand{\A}{\mathbb{A}}
\newcommand{\PP}{\mathbb{P}}
\newcommand{\E}{\mathbb{E}}
\newcommand{\eps}{\varepsilon}
\newcommand{\dd}{\,\mathrm{d}}
\begin{document}
\section*{UnsolvedMath ID: 30006595}
\textbf{Source problem code:} catalogue entry
\subsection{Definitions and mathematical statement}
The supplied record refers to integrals \(R_{k,\ell}(X)\), but does not supply their defining integrands, integration domains, normalization, admissible indices, or the domain of \(X\). For the depth-one family, the displayed expansion records the indexing
\[
\mathcal G(t,X)=\sum_{k=2}^{\infty}R_{k,1}(X)t^{k-2}.
\]
Here \(\mathcal G\) denotes the generating expression mentioned in the source; this notation does not assert convergence. At the level of formal coefficient indexing, its coefficient of \(t^m\) is \(R_{m+2,1}(X)\) for integers \(m\geq 0\). The supplied transcription of the hypergeometric identity labelled Theorem 2 is corrupted and does not determine its complete right-hand side reliably.

\textbf{Research question, subject to unresolved source definitions.} Determine whether the coefficients of higher powers of \(t\) in \(\mathcal G(t,X)\) admit interpretations as modular forms or higher-order modular antiderivatives. Identify and prove the corresponding structure for the integrals \(R_{k,\ell}(X)\) at depths \(\ell>1\). The original excerpt also raises the task of writing linear differential equations for the higher-power coefficients.

A precise modularity assertion requires the source definitions of the integrals, the admissible indices and \(X\)-domain, the relation between \(X\) and a modular parameter, the relevant modular group and transformation data, and the intended meaning of a higher-order modular antiderivative. The cutoff meant by ``higher powers'' is unspecified. The differential-equation variable, permitted coefficients, and desired order are also unspecified. Consequently, the supplied evidence supports this qualified research question but does not determine a fully quantified conjecture.

The supplied assessment says that modular identifications are available only for low powers and that higher powers and the general higher-depth structure remain difficult and unproved. The original excerpt specifically reports partial results for \(R_{k,2}(X)\). These are source-reported qualifications, not mathematical or literature conclusions independently verified here.
\subsection{Short English statement}
The question asks whether the higher coefficients of a generating function for \(R_{k,1}(X)\) have modular interpretations, and what analogous structure exists for \(R_{k,\ell}(X)\) with \(\ell>1\). The integral definitions and modularity conventions are missing, and the displayed hypergeometric formula is corrupted, so a fully precise conjecture requires source clarification.
\subsection{Sources}
\begin{itemize}
\item Local UnsolvedMath registry, permanent ID \(30006595\), source problem code \texttt{catalogue entry}, title ``Modularity of Higher-\(t\) Generating-Function Coefficients.'' The supplied record belongs to the Oberwolfach Reports Open Problems collection; no individual proposer is specified.
\item \emph{Multiple Zeta Values and Modular Forms} (2026), Oberwolfach Reports, DOI \texttt{10.4171/owr/2026/7}, pp.\ 433--486. Supplied source and literature-source link, labelled ``Oberwolfach Report 7/2026'': \url{https://doi.org/10.4171/owr/2026/7}. The background describes these reports as documenting mathematical workshops and their open-problem discussions.
\item Literature assessment dated \texttt{2026-08-22}: the supplied curation labels the problem open and reports modular identifications only for low powers, with higher powers and the general higher-depth structure remaining difficult and unproved. This is supplied dated curation, not an independently verified assessment; time-sensitive status and cited claims require independent verification.
\item The registry labels the statement \texttt{corrected\_verified} and states that it was reconstructed and checked against the official Oberwolfach report. This is a supplied verification claim; the report has not been independently consulted for this draft.
\end{itemize}
\end{document}
